\documentclass[11pt,a4paper]{article}
\usepackage[margin=25mm]{geometry}
\usepackage{amsmath,amssymb,amsthm,booktabs,longtable,hyperref}
\newtheorem{theorem}{Theorem}
\newtheorem{lemma}[theorem]{Lemma}
\newtheorem{proposition}[theorem]{Proposition}
\newcommand{\Aut}{\operatorname{Aut}}
\newcommand{\Cor}{\operatorname{Cor}}
\newcommand{\IP}{\mathrm{IP}}
\newcommand{\WIP}{\mathrm{WIP}}
\title{Historical names and exact shadow decisions for sporadic geometries}
\author{An application of the retained complete geometry records}
\date{}
\begin{document}\maketitle
\begin{abstract}
The shadow algorithms of Jacobs and Leemans give properties of actual
coset geometries. Their tables also use names from earlier classifications.
These are separate assertions: a list position or a matching answer does
not identify a historical geometry. We recover the available names from
the original authors' diagram supplements, using type indices, panel
orders, rank-two residues and correlations. Complete compatible families,
rather than arbitrarily selected representatives, suffice for the property
assertions. We establish 229 of the 232 positive occurrences through these profiles.
For the remaining three occurrences the complete $J_1$ classification is
stated by explicit mathematical profiles, while the historical ordinal
4.1 remains unverified. This completes a corrected principal-result
contract without claiming literal equivalence of historical table names.
\end{abstract}

\section{The mathematical and documentary inputs}
We use the complete sixteen-page edition of \cite{JL}, the common paper
\emph{Witnesses and orbits in finite incidence geometries}, and its
application \emph{Shadow intersections in sporadic and tree geometries}.
The latter supplies the full residue calculations, exact family maps,
source filters and corrected aggregate counts. Their written proofs and
complete computational records remain mathematical inputs here. In
particular, this note does not rerun or replace a subgroup census.

The retained population for each sporadic group consists of all firm,
residually connected, flag-transitive coset geometries in the stated
classification satisfying residual weak primitivity, local
2-transitivity and $\IP_2$. An equivalence is an automorphism of the
acting group on the whole family, followed by a permutation of types.
The complete earlier subgroup censuses and their accepted geometry tests
are retained for $M_{11},M_{12},J_1,J_2,M_{22},M_{23}$. For HS and $J_3$,
completeness of the historical population remains the classification
cited in \cite{JL}; the current finite certificates verify the entire
literal archive, its actual family quotient and all shadow decisions.
The distinction is essential: neither a diagram nor agreement of counts
proves that a finite list is complete.

The new documentary inputs are the diagram lists linked from the
original authors' publication pages. The archived PostScript files have
been retained in full, with retrieval URLs, timestamps and byte hashes.
Their PDF conversions are reading copies, not publisher editions.
Numbers and diagram symbols used below were read on the rendered pages;
the auxiliary byte-reading transcription is not a literal Unicode
transcription and is not a measurement coordinate. The complete
transcribed invariant list is \texttt{transcribed\_profiles.json}.
It does not contain any computed IP answer used to choose a family.

\begin{lemma}[Complete compatible families]\label{fibre}
Let $\mathcal C$ be a complete set of representatives for the permitted
$G$-geometries, and let a named source geometry $\Gamma$ belong to that
population. Suppose $\nu(\Gamma)=v$ is established independently of a
property $P$, where $\nu$ is invariant under the permitted equivalences.
Set $\mathcal C_v=\{C\in\mathcal C:\nu(C)=v\}$. If every element of
$\mathcal C_v$ has $P$, then $\Gamma$ has $P$. If $\mathcal C_v$ is a
singleton, it also identifies the source geometry up to the stated
equivalence. Without this uniqueness, the property conclusion does not
identify a particular marked representative.
\end{lemma}
\begin{proof}
Completeness gives a representative $C$ equivalent to $\Gamma$.
Invariance gives $C\in\mathcal C_v$. The hypothesis and invariance of
$P$ now give $P(\Gamma)$. A singleton fixes the equivalence class;
if the set is larger, the argument gives no choice between its members.
\end{proof}

For a family $(H_i)_{i\in I}$ the type sizes are $[G:H_i]$, and the
Borel subgroup is $B=\bigcap_iH_i$. A panel of type $i$ has size
$[\bigcap_{j\ne i}H_j:B]$. Thus the source diagram's order at that
node is this index minus one. We compare the complete multiset of
type-size/panel-order pairs when panel orders are needed, including
repeated type sizes; no dictionary overwrite of equal sizes is allowed.
For a source rank-two residue on types $i,j$, its incidence graph is
the coset graph in $K=\bigcap_{h\notin\{i,j\}}H_h$ with parabolics
$K\cap H_i,K\cap H_j$. This is the accepted residue theorem, not a
new inference from a diagram.

\section{Two correlation tests}\label{correlations}
All cosets in this section are right cosets, and $H^g=g^{-1}Hg$.
A correlation may permute types; an automorphism fixes types. The
type-preserving automorphism group of each geometry to which we apply
the following lemma is the stated faithful simple group $G$. This is
an explicitly retained fact of the corresponding source diagram list.

\begin{lemma}[Correlations normalizing the given action]\label{normalizing}
Suppose $G$ is flag-transitive on the coset geometry $(H_i)_{i\in I}$.
A correlation normalizing $G$ and inducing the type permutation $\pi$
exists if and only if an automorphism $\alpha\in\Aut(G)$ sends
$H_i$ to $H_{\pi(i)}$ for every $i$, after a simultaneous inner
adjustment. In particular, once the image chamber is returned to the
base chamber, the test is an exact whole-family automorphism test.
\end{lemma}
\begin{proof}
A normalizing correlation conjugates the faithful action of $G$ to
itself. Compose the correlation with a member of $G$ returning its
image of the base chamber to that chamber; flag transitivity permits
this. The induced automorphism of $G$ then sends each base-object
stabilizer to the stabilizer of the appropriate base object. Conversely,
such an automorphism defines $H_ix\mapsto H_{\pi(i)}\alpha(x)$.
The subgroup condition makes this well defined and bijective; since
$\alpha$ preserves intersections of cosets, it preserves incidence.
The initial inner adjustment is induced by a member of $G$ and does
not change existence.
\end{proof}

\begin{lemma}[Centralizing type interchange]\label{centralizing}
Let $\pi$ interchange two types and fix every other type. If $g\in G$
satisfies $H_i^g=H_{\pi(i)}$ for all $i$, then
\[
       H_ix\longmapsto H_{\pi(i)}g^{-1}x                         \tag{1}
\]
is a correlation commuting with the right action of $G$. Its square
is the identity if $g^2\in B$. Conversely, any correlation commuting
with $G$ and inducing $\pi$ is given by (1) for some such $g$, when
$G$ is flag-transitive.
\end{lemma}
\begin{proof}
The identity $g^{-1}H_i=H_{\pi(i)}g^{-1}$ proves that (1) is well
defined. It maps all object cosets by the same left multiplication,
so preserves all intersections, and commutes with every right
multiplication. Its square maps $H_ix$ to $H_ig^{-2}x$, which is the
same object for all $i,x$ when $g^2\in B$.

For the converse, the image of the base chamber has the form
$(H_{\pi(i)}y)_i$ for one $y\in G$ by flag transitivity.
Equivariance identifies the stabilizer of $H_i$ with the stabilizer
of $H_{\pi(i)}y$, giving $H_i=y^{-1}H_{\pi(i)}y$.
Put $g=y^{-1}$. Equivariance on every orbit gives (1).
\end{proof}

\paragraph{The two $M_{22}$ rank-three geometries.}
Source geometries 3.24 and 3.25 both have type sizes
$(672,9240,9240)$ and $|B|=4$. The two corresponding accepted families
are rank-three entries 3 and 5 in the complete prior $M_{22}$ record.
Their source lists distinguish them by the correlation group: 3.24
has no type interchange, whereas 3.25 has the outer extension.
Since their type-preserving automorphism group is $G=M_{22}$,
every correlation normalizes this action. Lemma~\ref{normalizing}
therefore gives a complete test in $\Aut(G)$.

The program first fixes the unique parabolic of order 660. Its
stabilizer in $\Aut(G)$ has order 1320. It finds an automorphism
carrying the first parabolic of order 48 to the second, if one exists,
and then searches the stabilizer of the second for an adjustment
carrying the other image back to the first. This searches a complete
transporter coset, not just one arbitrary representative. The test is
false for entry 3 and true for entry 5. The two simultaneous subgroup
equalities are checked when the test succeeds. Thus 3.24 is entry 3
and 3.25 is entry 5. Their already accepted shadow decisions give
IP for the former and WIP without IP for the latter.

\paragraph{The two $J_3$ diagram types.}
Source geometries 3.8 and 3.9 have the common type sizes
$(523260,523260,20520)$ and $|B|=8$. The possible actual inputs are
9,10,11 in the final literal assignment of \texttt{j3.rk3}.
The source gives $\Cor(3.8)=J_3\times2$ and
$\Cor(3.9)=\Aut(J_3)$, with $\Aut(\Gamma)=J_3$ in both cases.
The former centralizes $G$ by a type interchange; the latter does not,
because $G$ is centreless and $C_{\Aut(G)}(G)=1$.

In each actual family, the parabolic of order 2448 is self-normalizing
in $G$. An inner family interchange must normalize this parabolic,
so it is enough to enumerate its 2448 elements. There are respectively
8,0,8 elements interchanging the two parabolics of order 96 in inputs
9,10,11. In the positive cases each square lies in $B$.
Lemma~\ref{centralizing} identifies inputs 9 and 11 with the two
possibilities named 3.8, and input 10 with 3.9. The fact that 3.8
labels two geometries is expressly retained from its diagram list.
Both compatible inputs fail WIP, while input 10 has IP. No property
answer was used in either interchange test.

\section{The Mathieu residue discriminator}
The diagrams 4.7 and 4.8 have type sizes $(22,1540,176,330)$
and $|B|=8$. These give exactly rank-four entries 11 and 12 of the
complete prior $M_{22}$ record. The edge between the nodes of types
1540 and 330 is labelled 4 in 4.7 and has the ordinary label 3 in
4.8. Thus the corresponding rank-two residue graphs have girths 8
and 6, respectively.

For each candidate, form the subgroup $K$ fixing the other two types.
Construct every right coset of its two residue parabolics, with edges
defined by nonempty intersection. The complete small graphs have
\[
\begin{array}{c|rrr}
\text{prior entry}&|K|&\text{vertices}&\text{girth}\\\hline
11&360&30&8\\
12&168&14&6
\end{array}
\]
Both graphs are cubic and connected. Breadth-first distances from
every vertex determine their girths; the output includes every
adjacency list, and a separate Python traversal rechecks these
graph conclusions. Hence 4.7 is entry 11 and 4.8 is entry 12.
This calculation has no subgroup-census or shadow-search component.

Two other $M_{22}$ labels, 3.8 and 3.11, share type sizes
$(176,1540,176)$ and $|B|=18$. Their node orders are respectively
$(3,3,3)$ and $(3,1,3)$. The complete panel-index comparison
distinguishes the two actual families. In contrast, the two diagrams
4.15 and 4.16 are not separated by the invariants needed here. The
whole compatible set is retained, and every claimed shadow property
has the same value on both members. Lemma~\ref{fibre}, rather than
an invented choice of name, gives their property assertions.

\section{Transcription qualifications and source corrections}
Several numerical inconsistencies in auxiliary diagrams require care.
The $M_{12}$ diagram 3.4 prints a type index 432 alongside a
parabolic of order 192. Its index is $95040/192=495$, also printed
in the explicitly stated extension 4.9. The identifying index here
is 495, with the literal 432 preserved separately. Similarly,
$M_{23}$ diagram 3.9 prints 26585, but its parabolic of order 384
has index 26565, also printed in extension 4.11. Both diagram
contexts have been read. These corrections use group orders and
explicit source relations, not the desired IP outcome.

The auxiliary Borel descriptions at HS 3.4 and $M_{22}$ 3.7
conflict with the actual Borel orders in their uniquely matching
families. Those Borel strings are not used in the identification:
their type indices alone give unique retained classes. This note
does not claim a recovery of every group-structure assertion in
the auxiliary classifications. Its target is the named shadow
claims of \cite{JL}; the precise retained classification and group
identification inputs remain visible.

The archive README's HS 4.1 example is corrected by the actual
primary diagram. Its type sizes $(100,3850,3850,100)$ and Borel
order 192 select global class 31, literal input 45 (rank-four
active entry 1). The README instead points to active entry 2,
global class 32, whose certificate fails WIP. Thus the source's
positive geometry and the archive's illustrative pointer must
be distinguished. The earlier accepted HS 3.2 and $M_{12}$ 5.1
identifications are included in the same complete invariant ledger.
The separate HS 6.1 negative assertion is already resolved by the
unique retained rank-six class; it is outside the 232 positive
occurrences.

The corrected aggregate population remains that of the retained
application: $M_{12}$ rank three has split $(7,1,9)$ rather than
$(7,1,10)$, and the literal $J_3$ quotient has $(6,8,5)$ rather
than $(5,8,4)$. These triples list IP, WIP without IP, and failure
of WIP. No newly matched name changes an old shadow value or
creates a new geometry. In the $J_3$ supplement, labels 3.7 and
3.8 explicitly have two possibilities; their compatible actual
families account for the additional IP and non-WIP classes.

\section{Exact present boundary}
The accompanying table records every available source profile and
its entire compatible family. The machine-readable reconciliation
then treats every one of the 232 occurrences independently, at its
printed rank and with its printed IP or WIP-without-IP assertion.
All 229 occurrences supplied with documentary profiles agree on
every compatible representative. For the three occurrences naming $J_1$4.1, the corrected result is
instead the following complete mathematical classification. There is
one rank-three class, of type sizes $\{266,1463,1463\}$ and Borel
order4, which fails WIP. There are exactly two rank-four classes:
\[
\begin{array}{c|c|c|c}
\text{type sizes}&|B|&\IP_3,\IP_4&\WIP_3,\WIP_4\\\hline
\{266,1463,1463,2926\}&2&(1,1)&(1,1)\\
\{266,266,1463,2926\}&2&(0,0)&(0,0).
\end{array}
\]
These are the exact accepted complete census classes transported by
whole-family maps to the supplied inputs. They recover all three
aggregate $J_1$ counts in the source. The historical ordinal is not
needed in their mathematical statements.

The 2017 construction of Buekenhout--Connor identifies the required
type-size profile as 4.1 in Leemans's 2008 atlas and attributes its
earlier uniqueness to Gottschalk--Leemans. This establishes a later
name but, by itself, does not prove the same numbering in the 1998
chapter cited by \cite{JL}. We do not infer historical naming from
an archive position or from the one IP answer. The corrected principal contract
therefore uses the explicit profiles above. The three ordinal
correspondences remain unverified historical metadata, not proved
identities or new mathematical premises.

The complete required new computation is
\texttt{label\_invariants.g}, together with the entire accepted
$M_{22}$ and $J_3$ permutation input files, its full output and the
parsing/reconciliation programs. The large $J_3$ file is not priced
as a short selected witness. Old complete proof programs, data and
checkpoints remain acquired through the prior whole-file inventory.
Operational rendering and integrity checks are recorded separately.
No new formalization, conjecture answer, or numerical factor is
claimed. A later comparison must also treat the source's concrete
Magma listings symmetrically.

\small
\begin{longtable}{@{}llrlrl@{}}
\caption{Source invariant profiles and compatible current entries.
The entry numbers are internal record indices, not source suffixes.}\\
\toprule Group&Name&Page&Type sizes&$|B|$&Entries\\\midrule
\endfirsthead
\toprule Group&Name&Page&Type sizes&$|B|$&Entries\\\midrule
\endhead
\bottomrule\endfoot
$\mathrm{HS}$ & 3.1 & 4 & 176,15400,176 & 720 & 1 \\
$\mathrm{HS}$ & 3.2 & 4 & 176,15400,1100 & 720 & 2 \\
$\mathrm{HS}$ & 3.4 & 4 & 3850,28875,5775 & -- & 6 \\
$\mathrm{HS}$ & 3.7 & 5 & 4125,57750,1100 & 192 & 9 \\
$\mathrm{HS}$ & 3.10 & 5 & 1100,57750,1100 & 192 & 12 \\
$\mathrm{HS}$ & 3.13 & 6 & 352,30800,5600 & 72 & 16 \\
$\mathrm{HS}$ & 3.14 & 6 & 5600,154000,1100 & 72 & 18 \\
$\mathrm{HS}$ & 3.19 & 7 & 15400,61600,36960 & 40 & 28 \\
$\mathrm{HS}$ & 3.20 & 7 & 15400,346500,15400 & 32 & 29 \\
$\mathrm{HS}$ & 3.24 & 8 & 5600,30800,231000 & 16 & 34 \\
$\mathrm{HS}$ & 3.25 & 8 & 1100,154000,30800 & 16 & 36 \\
$\mathrm{HS}$ & 3.28 & 9 & 5600,184800,123200 & 12 & 39 \\
$\mathrm{HS}$ & 3.29 & 9 & 5600,231000,77000 & 12 & 41 \\
$\mathrm{HS}$ & 4.1 & 10 & 100,3850,3850,100 & 192 & 1 \\
$\mathrm{HS}$ & 4.3 & 10 & 1100,57750,4125,3850 & 64 & 3 \\
$\mathrm{HS}$ & 4.5 & 11 & 2200,4125,4125,3850 & 64 & 5 \\
$\mathrm{HS}$ & 4.7 & 11 & 1100,57750,57750,1100 & 32 & 7 \\
$\mathrm{HS}$ & 4.8 & 11 & 176,15400,61600,36960 & 20 & 8 \\
$\mathrm{HS}$ & 4.14 & 12 & 352,30800,5600,231000 & 8 & 18 \\
$\mathrm{HS}$ & 4.15 & 13 & 5600,1100,231000,5600 & 8 & 20 \\
$\mathrm{HS}$ & 4.16 & 13 & 352,231000,1100,352 & 8 & 22 \\
$\mathrm{HS}$ & 4.40 & 19 & 5600,462000,23100,5600 & 4 & 66 \\
$\mathrm{HS}$ & 5.6 & 21 & 5600,17600,123200,17600,5600 & 3 & 8 \\
$J_{2}$ & 3.1 & 4 & 560,840,525 & -- & 19 \\
$J_{2}$ & 3.2 & 4 & 100,3150,1575 & -- & 18 \\
$J_{2}$ & 3.3 & 4 & 1008,840,280 & -- & 17 \\
$J_{2}$ & 3.4 & 4 & 100,3150,280 & -- & 16 \\
$J_{2}$ & 3.7 & 5 & 525,3150,280 & -- & 15 \\
$J_{2}$ & 3.8 & 5 & 12600,1800,1800 & -- & 11 \\
$J_{2}$ & 3.9 & 5 & 1800,18900,6300 & -- & 12 \\
$J_{2}$ & 3.10 & 5 & 1800,25200,1800 & -- & 10 \\
$J_{2}$ & 3.16 & 6 & 4032,50400,10080 & -- & 4 \\
$J_{2}$ & 3.17 & 7 & 50400,10080,10080 & 2 & 3 \\
$J_{2}$ & 3.18 & 7 & 10080,50400,30240 & 2 & 1 \\
$J_{2}$ & 4.1 & 7 & 10080,4032,10080,4032 & 1 & 2 \\
$J_{2}$ & 4.2 & 8 & 4032,50400,10080,10080 & 1 & 1 \\
$J_{2}$ & 4.3 & 8 & 1575,100,100,1575 & 8 & 3 \\
$M_{11}$ & 3.3 & 4 & 66,220,22 & 6 & 1 \\
$M_{11}$ & 3.4 & 4 & 66,110,165 & 4 & 2 \\
$M_{11}$ & 3.5 & 4 & 66,220,165 & 4 & 3 \\
$M_{11}$ & 3.11 & 6 & 165,330,110 & 4 & 9 \\
$M_{11}$ & 3.12 & 6 & 165,330,66 & 4 & 10 \\
$M_{11}$ & 3.13 & 6 & 110,330,66 & 4 & 11 \\
$M_{11}$ & 3.18 & 7 & 11,55,165 & 8 & 13 \\
$M_{11}$ & 4.1 & 7 & 12,66,220,22 & 3 & 1 \\
$M_{11}$ & 4.4 & 8 & 12,66,220,165 & 2 & 3 \\
$M_{12}$ & 3.2 & 6 & 12,66,220 & 72 & 2 \\
$M_{12}$ & 3.4 & 6 & 66,220,495 & 16 & 4 \\
$M_{12}$ & 3.5 & 7 & 495,495,132 & 16 & 5 \\
$M_{12}$ & 3.6 & 7 & 66,1485,66 & 16 & 6 \\
$M_{12}$ & 3.7 & 7 & 66,1485,495 & 16 & 7 \\
$M_{12}$ & 3.9 & 7 & 220,495,132 & 12 & 9 \\
$M_{12}$ & 3.16 & 9 & 396,2970,396 & 8 & 16 \\
$M_{12}$ & 3.19 & 9 & 144,3960,1980 & 4 & 19 \\
$M_{12}$ & 4.9 & 12 & 12,66,220,495 & 8 & 9 \\
$M_{12}$ & 4.12 & 13 & 12,132,880,132 & 6 & 12 \\
$M_{12}$ & 4.19 & 15 & 144,3960,144,144 & 3 & 19 \\
$M_{12}$ & 4.20 & 15 & 1980,3960,144,144 & 2 & 20 \\
$M_{12}$ & 5.1 & 15 & 12,132,880,132,12 & 3 & 1 \\
$M_{22}$ & 3.1 & 4 & 22,231,77 & 192 & 25 \\
$M_{22}$ & 3.3 & 4 & 330,1155,231 & 64 & 24 \\
$M_{22}$ & 3.6 & 5 & 231,1540,176 & 24 & 21 \\
$M_{22}$ & 3.7 & 5 & 176,330,1540 & -- & 20 \\
$M_{22}$ & 3.8 & 5 & 176,1540,176 & 18 & 17 \\
$M_{22}$ & 3.11 & 6 & 176,1540,176 & 18 & 13 \\
$M_{22}$ & 3.12 & 6 & 770,6160,176 & 18 & 16 \\
$M_{22}$ & 3.13 & 6 & 176,6160,176 & 18 & 14 \\
$M_{22}$ & 3.17 & 7 & 616,13860,616 & 8 & 12 \\
$M_{22}$ & 3.18 & 7 & 616,3696,3696 & 4 & 7 \\
$M_{22}$ & 3.23 & 8 & 672,18480,3696 & 4 & 2 \\
$M_{22}$ & 3.24 & 8 & 672,9240,9240 & 4 & 3 \\
$M_{22}$ & 3.25 & 8 & 672,9240,9240 & 4 & 5 \\
$M_{22}$ & 4.3 & 9 & 22,231,1540,176 & 12 & 13 \\
$M_{22}$ & 4.7 & 10 & 22,1540,176,330 & 8 & 11 \\
$M_{22}$ & 4.8 & 10 & 22,1540,176,330 & 8 & 12 \\
$M_{22}$ & 4.9 & 10 & 77,2310,2310,77 & 8 & 8 \\
$M_{22}$ & 4.11 & 11 & 176,1540,231,176 & 6 & 6 \\
$M_{22}$ & 4.12 & 11 & 462,672,672,462 & 3 & 5 \\
$M_{22}$ & 4.13 & 11 & 462,672,672,462 & 3 & 5 \\
$M_{22}$ & 4.14 & 11 & 1232,672,6160,1232 & 3 & 4 \\
$M_{22}$ & 4.15 & 12 & 462,3696,672,9240 & 2 & 1, 2 \\
$M_{22}$ & 4.16 & 12 & 462,3696,672,9240 & 2 & 1, 2 \\
$M_{22}$ & 5.2 & 12 & 22,231,1540,176,176 & 3 & 6 \\
$M_{23}$ & 3.2 & 3 & 23,253,1771 & 960 & 13 \\
$M_{23}$ & 3.3 & 4 & 253,1771,253 & 384 & 11 \\
$M_{23}$ & 3.5 & 4 & 506,3795,1771 & 192 & 9 \\
$M_{23}$ & 3.6 & 4 & 506,8855,1771 & 72 & 8 \\
$M_{23}$ & 3.9 & 5 & 253,26565,3795 & 48 & 5 \\
$M_{23}$ & 3.13 & 6 & 1288,106260,28336 & 12 & 2 \\
$M_{23}$ & 4.2 & 6 & 23,253,1771,253 & 192 & 20 \\
$M_{23}$ & 4.4 & 6 & 506,3795,1771,23 & 64 & 19 \\
$M_{23}$ & 4.8 & 7 & 253,1771,8855,506 & 24 & 15 \\
$M_{23}$ & 4.11 & 8 & 253,26565,3795,7590 & 16 & 10 \\
$M_{23}$ & 4.12 & 8 & 253,26565,506,7590 & 16 & 11 \\
$M_{23}$ & 4.17 & 10 & 1771,1288,28336,14168 & 6 & 5 \\
$M_{23}$ & 4.19 & 10 & 1288,212520,5313,14168 & 4 & 3 \\
$M_{23}$ & 5.3 & 11 & 23,253,1771,8855,506 & 12 & 5 \\
$M_{23}$ & 5.8 & 13 & 23,253,8855,253,506 & 8 & 4 \\
$M_{23}$ & 5.9 & 13 & 1771,1288,28336,14168,23 & 3 & 2 \\
$J_{3}$ & 3.1 & 4 & 46512,43605,17442 & 36 & 1 \\
$J_{3}$ & 3.2 & 4 & 23256,17442,6156 & 30 & 2 \\
$J_{3}$ & 3.3 & 4 & 26163,261630,6156 & 24 & 3 \\
$J_{3}$ & 3.4 & 4 & 23256,261630,43605 & 24 & 4 \\
$J_{3}$ & 3.5 & 4 & 43605,43605,46512 & 24 & 5 \\
$J_{3}$ & 3.6 & 5 & 14688,14688,837216 & 10 & 6 \\
$J_{3}$ & 3.7 & 5 & 20520,1569780,523260 & 8 & 7, 8 \\
$J_{3}$ & 3.8 & 5 & 523260,523260,20520 & 8 & 9, 11 \\
$J_{3}$ & 3.9 & 5 & 523260,523260,20520 & 8 & 10 \\
$J_{3}$ & 3.10 & 5 & 20520,1569780,69768 & 8 & 12 \\
$J_{3}$ & 3.11 & 6 & 20520,20520,69768 & 8 & 13 \\
$J_{3}$ & 3.12 & 6 & 20520,20520,139536 & 8 & 15 \\
$J_{3}$ & 3.13 & 6 & 20520,69768,1046520 & 8 & 16 \\
$J_{3}$ & 3.14 & 6 & 20520,20520,1046520 & 8 & 18 \\
$J_{3}$ & 3.15 & 6 & 20520,930240,139536 & 6 & 19 \\
$J_{3}$ & 3.16 & 7 & 20520,20520,523260 & 6 & 20 \\
$J_{3}$ & 3.17 & 7 & 20520,930240,20520 & 6 & 21 \\
\end{longtable}

\normalsize
\begin{thebibliography}{9}
\bibitem{JL} P. Jacobs and D. Leemans,
An algorithmic analysis of the intersection property,
\emph{LMS J. Comput. Math.} \textbf{7} (2004), 284--299.
\bibitem{supp} M. Dehon, D. Leemans and X. Miller, $M_{11}$ diagram
list; M. Dehon and D. Leemans, supplement on $M_{12}$; D. Leemans,
supplements on $J_2,J_3,M_{22},M_{23}$ and HS.
Exact author-hosted archived editions, conversion hashes and
physical diagram pages are listed in the documentary source ledger.
\bibitem{BC} F. Buekenhout and T. Connor,
Alternative existence proof of the geometry of Ivanov--Shpectorov,
\emph{Innov. Incidence Geom.} \textbf{15} (2017), 73--121,
especially Section 4.3 and Figure 8.
\end{thebibliography}
\end{document}
